% !TEX root = ../main.tex

This chapter reviews the main strands of literature on which the thesis is grounded and shows how they converge within the proposed methodological framework. Rather than offering a generic overview of related work, it reconstructs the conceptual trajectory that motivates the thesis. The discussion begins with Explainable Artificial Intelligence (XAI), which provides the broader interpretive context in which this work is situated. It then turns to counterfactual explanations, which constitute the central explanatory perspective adopted in the thesis. From there, the chapter moves to feature selection and association-rule mining, two areas treated here not as isolated topics, but as essential components for the structured aggregation and analysis of counterfactual evidence.

This organization directly reflects the specific contribution of the thesis. The aim of the review is not merely to catalogue prior methods, but to show why the combination of local counterfactual evidence, feature-level relevance, and pattern-level synthesis helps address a gap that remains only partially covered in the existing literature. Accordingly, the chapter alternates between general survey material and explicit thesis positioning, so that each section contributes both to the state of the art and to the rationale underlying the proposed approach.

\section{Explainable Artificial Intelligence (XAI)}
Explainable Artificial Intelligence (XAI) is a field of research devoted to making machine-learning models and their predictions intelligible to human users. According to \textcite{adadi2018peeking}, \textcite{arrieta2020explainable}, and \textcite{molnar2022iml}, XAI seeks to clarify which factors influence a decision and how they contribute to the final output. The need for such methods has grown rapidly with the adoption of high-performing yet inherently opaque models, especially complex ensembles and deep neural networks, whose predictive accuracy is often achieved at the expense of interpretability. In the survey literature, \textcite{adadi2018peeking} and \textcite{arrieta2020explainable} describe this trade-off between predictive power and transparency as one of the central challenges of modern AI.

Within this context, XAI aims to make model behavior understandable and verifiable, enabling practitioners to validate learned patterns and detect potential failure modes. Explanations also support human decision-makers in high-stakes settings, where automated predictions must be interpreted critically rather than accepted passively. More broadly, explainability serves an important governance function by supporting accountability, trust, fairness assessment, and compliance-oriented reporting. As emphasized by \textcite{adadi2018peeking} and \textcite{arrieta2020explainable}, these needs arise across a wide range of application domains, although the concrete purpose of explanation varies with the decision context. In healthcare and biomedicine, explanations support clinical interpretation and make risk predictions more transparent. In finance and credit-risk analysis, they help justify automated assessments to analysts, auditors, and affected individuals. In industrial predictive-maintenance settings, they assist engineers in diagnosing failure patterns, distinguishing causal signals from noise, and prioritizing interventions. Similar demands arise in transportation and public-sector decision-support systems. For these reasons, XAI is now widely regarded as a foundational component of responsible AI pipelines rather than as a purely optional post-processing step.

\subsection{Bias, Fairness, and Their Connection with XAI}
The widespread deployment of machine-learning systems in socially relevant domains has made it clear that predictive performance alone is not sufficient to assess the quality of an AI system. Models often inherit and reproduce biases already present in historical data, data-collection procedures, labeling practices, or broader modeling choices. When these distortions disproportionately affect particular individuals or demographic groups, they raise serious fairness concerns. As discussed by Barocas and Selbst \cite{barocas2016big}, Mehrabi et al. \cite{mehrabi2021}, and, more broadly, Arrieta et al. \cite{arrieta2020explainable}, this issue is especially critical in domains such as credit scoring, hiring, insurance, healthcare, and public decision-support, where automated outputs can directly influence access to resources, opportunities, and rights.

In the machine-learning literature, \emph{bias} is a broad notion that refers to multiple sources of systematic distortion. A dataset may be biased because some populations are underrepresented, because labels reflect past discriminatory practices, or because some variables act as direct or indirect proxies for sensitive attributes. \emph{Fairness}, by contrast, concerns the extent to which a model and its decisions satisfy normative requirements of equitable treatment across individuals or groups. For this reason, fairness is typically studied through families of formal criteria, such as demographic parity, equal opportunity, equalized odds, calibration, and counterfactual fairness. As reviewed by Mehrabi et al. \cite{mehrabi2021}, discussed by Barocas and Selbst \cite{barocas2016big}, and extended in the counterfactual setting by Kusner et al. \cite{kusner2017counterfactual}, each criterion captures a different intuition of what a fair automated decision process should look like. Because these criteria are often mutually incompatible, fairness analysis is inherently contextual and requires explicit choices about which harms should be monitored and mitigated.

This close relationship is one of the main reasons fairness is so strongly connected to XAI. When a model produces systematically different outcomes across subpopulations, explanation methods are needed to determine whether such disparities arise from valid predictive structure, spurious correlations, or the undue influence of proxy variables. As emphasized by Arrieta et al. \cite{arrieta2020explainable} and Adadi and Berrada \cite{adadi2018peeking}, XAI methods support this diagnostic process by revealing which features drive predictions, highlighting whether local decisions depend on problematic variables, and helping assess whether the model relies on stable and plausible patterns rather than data artifacts. In this sense, explainability does not by itself guarantee fairness, but it does provide the analytical tools needed to inspect, contest, and document potentially unfair behavior.

The relationship is also bidirectional, because fairness requirements influence which explanations are considered useful in practice. In high-stakes settings, stakeholders want to know not only \emph{why} a prediction was made, but also whether similar individuals are treated consistently, whether sensitive information has implicitly influenced the output, and whether actionable changes can be suggested without reproducing discriminatory assumptions. From this perspective, XAI and fairness should be understood as complementary dimensions of responsible AI: fairness defines part of the evaluative objective, while XAI provides the interpretive mechanisms through which that objective can be examined and operationalized. This connection is particularly strong in the case of counterfactual explanations. Wachter et al. \cite{wachter2017counterfactual} show how such explanations reveal the role of modifiable attributes, Ustun et al. \cite{ustun2019actionable} highlight their relevance for actionable recourse, and Kusner et al. \cite{kusner2017counterfactual} connect them to asymmetric treatment patterns and fairness-sensitive forms of reasoning.

\subsection{XAI Methods: Classification and Taxonomy from Survey Literature}
Survey papers generally avoid forcing XAI methods into a single, mutually exclusive taxonomy. Instead, as shown by Guidotti et al. \cite{guidotti2018survey}, Adadi and Berrada \cite{adadi2018peeking}, and Arrieta et al. \cite{arrieta2020explainable}, they classify explanation methods along multiple complementary axes. In other words, XAI methods are not assigned to a rigid class, but are described through a combination of dimensions such as scope, access to the predictive model, and the form of the explanation.

A first major distinction separates \emph{intrinsic} from \emph{post-hoc} explainability. Intrinsic models, such as decision trees, sparse linear models, generalized additive models, and rule-based classifiers, are transparent by design and can therefore be inspected directly \cite{molnar2022iml,arrieta2020explainable}. Post-hoc methods, by contrast, are applied after training to explain the behavior of an otherwise opaque predictor. Representative examples include LIME \cite{ribeiro2016lime}, SHAP \cite{lundberg2017shap}, counterfactual explanations \cite{wachter2017counterfactual}, partial dependence plots \cite{friedman2001pdp}, and surrogate-rule extraction approaches \cite{guidotti2018survey,adadi2018peeking}.

A second major dimension concerns the scope of the explanation. Global explanations aim to reconstruct the model's overall behavior across a dataset. Classic examples include partial dependence analysis introduced by Friedman \cite{friedman2001pdp} and the broader global-explanation perspective reviewed by Arrieta et al. \cite{arrieta2020explainable}. Local explanations, by contrast, justify individual predictions and are especially useful when specific decisions must be defended. This local perspective is exemplified by LIME \cite{ribeiro2016lime}, SHAP \cite{lundberg2017shap}, counterfactual explanations \cite{wachter2017counterfactual}, and the wider taxonomy summarized by Guidotti et al. \cite{guidotti2018survey}.

Methods are also classified according to their relationship with the predictive model. \emph{Model-specific} approaches exploit architectural properties internal to the model, such as gradient-based saliency for neural networks or split-based importance measures for tree ensembles, a distinction emphasized by Arrieta et al. \cite{arrieta2020explainable}. \emph{Model-agnostic} approaches, instead, rely only on input--output behavior, which makes them applicable across heterogeneous predictive systems. Representative examples include LIME \cite{ribeiro2016lime}, SHAP \cite{lundberg2017shap}, and counterfactual explanations \cite{wachter2017counterfactual}; broader survey discussions by Guidotti et al. \cite{guidotti2018survey} and Adadi and Berrada \cite{adadi2018peeking} place these methods within the general model-agnostic family.

Finally, the literature distinguishes methods by the form of explanation they produce. Feature-attribution methods, such as SHAP \cite{lundberg2017shap} and saliency-based approaches discussed by Arrieta et al. \cite{arrieta2020explainable}, quantify the contribution of individual variables. Example-based methods instead use representative instances, including prototypes, that is, typical examples of a class; criticisms, that is, cases poorly represented by those prototypes; and counterfactuals, following the perspective of Wachter et al. \cite{wachter2017counterfactual} and Guidotti et al. \cite{guidotti2018survey}. Surrogate approaches approximate black-box behavior through simpler models, as in Ribeiro et al. \cite{ribeiro2016lime} and Molnar \cite{molnar2022iml}, whereas rule-based explanations express predictive logic through human-readable rules, a family discussed by both Guidotti et al. \cite{guidotti2018survey} and Arrieta et al. \cite{arrieta2020explainable}. Within this broad taxonomy, counterfactual explanation is the XAI method most central to this thesis because it provides the primary explanatory perspective through which feature relevance is analyzed.


\begin{table}[H]
\centering
\renewcommand{\arraystretch}{1.18}
\caption{Summary of the main taxonomy dimensions used to classify XAI methods.}
\label{tab:xai-taxonomy-summary}
\begin{tabular}{>{\raggedright\arraybackslash}p{0.22\textwidth} >{\raggedright\arraybackslash}p{0.24\textwidth} >{\raggedright\arraybackslash}p{0.39\textwidth}}
  \toprule
  \textbf{Dimension} & \textbf{Categories} & \textbf{Representative methods and models} \\
  \midrule
  Explainability design & Intrinsic; post-hoc & Decision trees, sparse linear models,\newline generalized additive models;\newline LIME, SHAP, counterfactual explanations,\newline partial dependence plots \\
  Explanation scope & Global; local & Global surrogate models,\newline partial dependence analysis,\newline aggregated feature-importance measures;\newline LIME, SHAP, counterfactual explanations \\
  Relation to predictive model & Model-specific; model-agnostic & Saliency-based methods for neural networks,\newline gain- or split-based importance measures for tree ensembles;\newline LIME, SHAP, permutation-based explanations,\newline counterfactual explanations \\
  Explanation form & Feature-attribution; example-based; surrogate; rule-based & SHAP and saliency-based methods;\newline prototypes, criticisms, and counterfactuals;\newline local linear surrogates and global interpretable proxies;\newline rule-extraction methods \\
  \bottomrule
\end{tabular}
\renewcommand{\arraystretch}{1.0}
\end{table}

\section{Counterfactual Explanations}
Counterfactual explanations answer explicit \enquote{what-if} questions by identifying the minimal changes to an input required to alter a model's prediction toward a desired outcome \cite{wachter2017counterfactual}. Intuitively, they describe how an instance must cross a decision boundary. In a lending scenario, for example, a counterfactual explanation can clarify why an application was rejected and which realistic changes could lead to approval. As highlighted by Chou et al. \cite{chou2022counterfactuals}, this contrastive structure is especially effective because it addresses the question of why one outcome occurred \emph{instead of} another plausible alternative. Their practical appeal lies in the fact that they express explanations within the same feature space as the original decision, thereby translating abstract model behavior into concrete feature-level variations.

Viewed through the taxonomy outlined above, the position of counterfactuals is relatively clear. Guidotti et al. \cite{guidotti2018survey} and Arrieta et al. \cite{arrieta2020explainable} generally classify them as post-hoc and local methods, while Wachter et al. \cite{wachter2017counterfactual} and Chou et al. \cite{chou2022counterfactuals} emphasize their model-agnostic and contrastive character. They are also typically interpreted as example-based explanations because they justify a decision by presenting a nearby alternative instance associated with a different outcome. This positioning is especially important for the present thesis: counterfactuals are not used here as global surrogates or as intrinsically interpretable models, but as local explanatory objects from which systematic information about feature relevance can be extracted. They occupy a particularly useful position within the XAI landscape because they preserve decision-specific detail while still allowing structured aggregation across many cases.

To be informative and practically useful, counterfactual explanations are typically evaluated against a set of core quality criteria. A well-constructed counterfactual should satisfy \emph{proximity} to the original instance while maintaining \emph{sparsity}, so that as few features as possible are changed and the explanation remains understandable. Moreover, the suggested changes cannot be purely mathematical artifacts: \emph{plausibility} requires the counterfactual to remain consistent with the underlying data distribution, while \emph{actionability} requires modified features to correspond to realistic and implementable interventions. Finally, when multiple recourse paths exist, \emph{diversity} becomes important because it gives users alternative ways of achieving a favorable outcome. Here, \emph{recourse} refers to the set of feasible changes through which an individual may obtain a desired decision outcome. This requirement is emphasized, from different perspectives, by Russell \cite{russell2019efficient}, Mothilal et al. \cite{mothilal2020explaining}, Poyiadzi et al. \cite{poyiadzi2020face}, and Karimi et al. \cite{karimi2022survey}.

These design principles make counterfactuals valuable not only for communicating individual predictions, but also for diagnosing local decision boundaries and identifying non-actionable dependencies, that is, cases in which the model relies on variables that the subject cannot realistically change. Counterfactuals are also useful for comparing alternative recourse paths and testing the coherence of a predictor under small perturbations, as shown in the foundational formulation of Wachter et al. \cite{wachter2017counterfactual}, the actionable-recourse perspective of Ustun et al. \cite{ustun2019actionable}, and the diversity-oriented framework of Mothilal et al. \cite{mothilal2020explaining}.

In the context of fairness, counterfactual reasoning plays a particularly important role. As explored by Kusner et al. \cite{kusner2017counterfactual} through the notion of \emph{counterfactual fairness}, these methods test whether a prediction remains stable when a sensitive attribute is hypothetically changed while all other relevant factors are held fixed. They are therefore highly informative for fairness auditing, because they can expose asymmetric recourse burdens, reveal when some demographic groups must undergo substantially larger changes to obtain a favorable outcome, and highlight decision patterns shaped by socially correlated variables. This function emerges clearly in the work of Kusner et al. \cite{kusner2017counterfactual}, Ustun et al. \cite{ustun2019actionable}, and Karimi et al. \cite{karimi2022survey}. Although counterfactuals do not replace formal fairness criteria, they provide a particularly concrete mechanism for inspecting unfair behavior.

For this thesis, the most important aspect of counterfactual explanations is their relationship with feature importance. In general, \emph{feature importance} refers to the extent to which an input variable contributes to a model's prediction, while \emph{relevant features} are those variables that carry decision-useful information for the task under study. Depending on the explanatory framework, relevance may refer to statistical contribution, predictive utility, or a more causal and decision-oriented notion of influence on the output. In the perspective adopted here, counterfactual analysis gives these notions an operational interpretation: a feature is considered especially relevant when modifying it is repeatedly necessary to flip the prediction, and its importance increases when it recurs across many valid and minimal counterfactuals. This yields a notion of \emph{decision-relevant feature importance} that is based not on variance decomposition alone, but on necessity for altering the prediction. In particular, Alfeo et al. \cite{alfeo2023bocsor} show that although counterfactuals are inherently local, aggregating them can yield stable global relevance signals. Closely related work by Mothilal et al. \cite{mothilal2021explaining} and Chou et al. \cite{chou2022counterfactuals} further shows that counterfactual explanations and feature-attribution methods should be viewed as complementary perspectives. This insight directly motivates the methodological approach of the thesis: local counterfactual structure is used as a bridge toward global, feature-level interpretation, and the resulting recurring patterns are subsequently summarized through association-based analysis.

\section{Feature Selection}
Feature selection aims to identify a subset of relevant variables that preserves or improves predictive utility while reducing redundancy and noise, as classically defined by Guyon and Elisseeff \cite{guyon2003introduction}. Beyond computational efficiency, it is also used to improve model interpretability, stability, and robustness. In the survey literature, Chandrashekar and Sahin \cite{chandrashekar2014survey} and Li et al. \cite{li2017feature} generally organize feature-selection methods into three broad families: \emph{filter} methods, which score features independently of the predictive model; \emph{wrapper} methods, which evaluate subsets through model performance; and \emph{embedded} methods, which integrate selection into the training process itself. Li et al. \cite{li2017feature} further extend this taxonomy to include distinctions such as supervised versus unsupervised selection, single-feature versus group selection, and optimization-based criteria.

Filter methods typically rely on mutual information, statistical tests, or correlation-based criteria. Wrapper methods search directly over subsets and often achieve strong predictive performance at a higher computational cost. Embedded methods include sparsity-inducing regularization and tree-based importance mechanisms, as discussed by Chandrashekar and Sahin \cite{chandrashekar2014survey} and by Li et al. \cite{li2017feature}. The same surveys also emphasize that no single family is uniformly superior; each embodies a different trade-off among predictive performance, computational cost, and interpretability.

Here, feature selection is approached through feature importance derived from counterfactual explanations. The underlying intuition is straightforward: if certain variables repeatedly appear in minimal counterfactual changes, they can be interpreted as especially relevant to crossing the decision boundary and revising the prediction. In this sense, counterfactual evidence does not merely accompany feature selection; it provides the explanatory basis on which variables are prioritized and constrained subsets are identified, especially when the objective is to emphasize actionable variables rather than simple predictive compression. This interpretation is consistent with the counterfactual relevance perspective developed by Alfeo et al. \cite{alfeo2023bocsor} and with the broader conceptual framing discussed by Chou et al. \cite{chou2022counterfactuals}. Recent work also investigates how instance-level counterfactual evidence can isolate features that are crucial for changing decisions in local neighborhoods, thereby complementing more classical feature-selection criteria, as discussed by Chen et al. \cite{chen2022causalfs} and Li et al. \cite{li2017feature}.

A related line of research explores association-rule-based feature selection. In this case, frequent rules identify variables that recur in discriminative combinations, so that importance emerges from interaction patterns rather than from marginal relevance alone. As noted by Hullermeier \cite{hullermeier2006association} and Ramírez-Gallego et al. \cite{ramirezgallego2014association}, rule-based methods are particularly useful when the goal is to preserve interpretability while still capturing multivariate structure. For the purposes of this thesis, this creates a natural complementarity: counterfactual explanations isolate the variables involved in a decision change, whereas association-oriented analysis summarizes how those variables interact across many explanations. Therefore, rather than proposing a new predictive feature-selection algorithm, this work uses counterfactual evidence to rank decision-relevant variables and to generate structured evidence that can support feature-selection reasoning.


\begin{table}[H]
\centering
\renewcommand{\arraystretch}{1.18}
\caption{Summary of the main taxonomy dimensions used to classify feature-selection methods.}
\label{tab:feature-selection-taxonomy-summary}
\begin{tabular}{>{\raggedright\arraybackslash}p{0.22\textwidth} >{\raggedright\arraybackslash}p{0.24\textwidth} >{\raggedright\arraybackslash}p{0.39\textwidth}}
  \toprule
  \textbf{Dimension} & \textbf{Categories} & \textbf{Representative criteria or methods} \\
  \midrule
  Selection family & Filter; wrapper; embedded & Mutual information, statistical tests; subset search with predictive models; L1 regularization, tree-based selection \\
  Learning setting & Supervised; unsupervised & Label-dependent relevance criteria; variance, similarity, or redundancy criteria \\
  Selection granularity & Single-feature; group-based & Ranking individual variables; selecting blocks or structured feature groups \\
  Optimization view & Ranking; subset search; regularized optimization & Score-based ordering; heuristic or exhaustive subset evaluation; sparsity-constrained training \\
  \bottomrule
\end{tabular}
\renewcommand{\arraystretch}{1.0}
\end{table}

\section{Association-Rule Mining: Taxonomy, Survey Perspective, and Relevance for This Thesis}
Association-rule mining (ARM) is a descriptive data-mining framework devoted to discovering recurring co-occurrence patterns. In its standard notation, a rule is written as $A \Rightarrow B$, where $A$ and $B$ are sets of items, usually called \emph{itemsets}. Intuitively, the rule states that when the items in $A$ appear together, the items in $B$ also tend to appear. In market-basket analysis, an item may literally be a purchased product; in more general analytical contexts, such as the one considered in this thesis, an item can instead represent a feature or a discretized feature-value condition. The classical foundations of ARM lie in the work of Agrawal and Srikant \cite{agrawal1994fast} and in the later pattern-growth perspective developed by Han et al. \cite{han2000mining}. Although originally introduced for market-basket analysis, ARM has since been applied to many other domains, including bioinformatics, fraud detection, and medical decision support, as discussed by Hullermeier \cite{hullermeier2006association} and Fournier-Viger et al. \cite{fournierviger2017survey}. Of particular relevance to this thesis is the work of Pedreschi, Ruggieri, and Turini \cite{pedreschi2008discriminationaware}, who proposed the use of association rules specifically for discrimination detection in automated decision systems. Their framework defines discriminatory classification rules as those whose confidence changes significantly when conditioned on a protected attribute, and it demonstrates that ARM can serve as a principled tool for identifying and quantifying bias in data-driven decisions. This line of research provides direct precedent for the interpretive use of association rules adopted in the present thesis, where extracted rules are examined for the presence of sensitive attributes as an indicator of potential bias in the classifier's decision logic. More broadly, the relevance of ARM lies in its ability to reveal structural interactions that often remain hidden in univariate analyses.

From a methodological standpoint, ARM proceeds in two main stages. First, it identifies recurrent feature combinations that satisfy a minimum \emph{support} threshold, that is, a minimum frequency of occurrence in the dataset. These combinations are usually called \emph{frequent itemsets}: an itemset is frequent when it appears in at least a user-defined proportion or count of transactions. Second, rules are derived from these frequent itemsets subject to additional interestingness constraints. In the rule $A \Rightarrow B$, $A$ is called the \emph{antecedent} and $B$ the \emph{consequent}. Among the most common evaluation measures are \emph{confidence}, which quantifies how reliably the consequent follows from the antecedent, and \emph{lift}, which compares the observed co-occurrence with what would be expected under statistical independence. These measures are central both in the classical presentations of Agrawal and Srikant \cite{agrawal1994fast} and Han et al. \cite{han2000mining}, and in later discussions of interestingness criteria such as that of McNicholas and Murphy \cite{mcnicholas2008lift}. As Hullermeier \cite{hullermeier2006association} emphasizes, they are essential for distinguishing trivial frequency from genuinely informative patterns.

The computation of frequent itemsets has historically been dominated by two major algorithmic paradigms. The first is candidate generation, exemplified by Apriori \cite{agrawal1994fast}, which builds larger itemsets level by level from smaller frequent ones and prunes the search space through the anti-monotonicity of support. The second is pattern growth, exemplified by FP-Growth \cite{han2000mining}, which avoids generating large numbers of explicit candidates by compressing the database into a prefix tree and recursively extracting frequent patterns from it. This distinction is especially important for the present thesis because the implementation adopted in the project relies specifically on FP-Growth. The choice is motivated by its practical efficiency when many recurring combinations must be explored and explicit candidate generation would become unnecessarily expensive.

Survey literature commonly organizes ARM along several complementary taxonomic dimensions. Hullermeier \cite{hullermeier2006association} and Fournier-Viger et al. \cite{fournierviger2017survey} both emphasize this multidimensional view. One dimension concerns the \emph{pattern type}, which ranges from classical frequent itemsets to sequential and high-utility patterns, moving from the foundational framework of Han et al. \cite{han2000mining} to the broader survey synthesis of Fournier-Viger et al. \cite{fournierviger2017survey}. A second dimension concerns the \emph{rule type}, which may be Boolean, quantitative, or fuzzy. Boolean rules express simple presence-or-absence conditions, quantitative rules involve numerical intervals or thresholds, and fuzzy rules allow gradual membership rather than sharp boundaries. This distinction is useful for positioning ARM within the broader literature; however, in the present thesis ARM is used only in its Boolean form, that is, by encoding whether an item is present or absent, and no fuzzy logic is adopted. A third dimension concerns the \emph{constraint setting}, including support thresholds, rule-length limits, and redundancy filters, all of which are essential because unconstrained mining often produces extremely large rule sets. A fourth dimension concerns the \emph{interestingness criterion}, which extends beyond support to include lift, leverage, and certainty factors. Finally, the \emph{algorithmic family} dimension contrasts candidate-generation methods such as Apriori \cite{agrawal1994fast} with pattern-growth approaches such as FP-Growth \cite{han2000mining}, directly reflecting the scalability issues discussed in the survey literature \cite{fournierviger2017survey}.

\begin{table}[H]
\centering
\renewcommand{\arraystretch}{1.18}
\caption{Summary of the main taxonomy dimensions used to classify association-rule mining methods.}
\label{tab:arm-taxonomy-summary}
\begin{tabular}{>{\raggedright\arraybackslash}p{0.22\textwidth} >{\raggedright\arraybackslash}p{0.24\textwidth} >{\raggedright\arraybackslash}p{0.39\textwidth}}
  \toprule
  \textbf{Dimension} & \textbf{Categories} & \textbf{Representative interpretation} \\
  \midrule
  Pattern type & Frequent; sequential; high-utility & Co-occurring itemsets; ordered patterns; quantity/value-aware patterns \\
  Rule type & Boolean; quantitative; fuzzy & Presence/absence rules; interval-based rules; gradual-membership rules \\
  Interestingness criteria & Support; confidence; lift & Frequency; conditional reliability; deviation from independence \\
  Algorithm family & Candidate generation; pattern growth & Apriori-style search; FP-Growth-style compressed mining \\
  Thesis setting & Boolean ARM; FP-Growth & Presence/absence encoding; efficient mining of frequent explanatory patterns \\
  \bottomrule
\end{tabular}
\renewcommand{\arraystretch}{1.0}
\end{table}

A crucial methodological issue concerns the treatment of multiplicities or quantities.
In classical ARM, each item is encoded in Boolean form: it is either present or absent in a transaction. This representation is not merely a simplifying assumption; it is the foundation that makes the anti-monotonicity of support operationally useful.
The anti-monotonicity property states that if an itemset is frequent, then all of its subsets must also be frequent; equivalently, if an itemset is infrequent, then all of its supersets must also be infrequent. This property allows classical frequent-itemset mining algorithms to efficiently prune large portions of the search space \cite{agrawal1994fast,han2000mining}.

When multiplicities are treated explicitly (e.g., distinguishing between one occurrence and multiple occurrences of the same item), this clean subset--superset ordering breaks, and the search space can no longer be reduced through the same anti-monotonic logic.
For this reason, multiplicities are not considered here: items are encoded strictly through presence or absence. This methodologically motivated choice allows the analysis to remain within the classical Boolean ARM framework, preserving the pruning structure required for efficient frequent-itemset mining, and aligns fully with the use of FP-Growth adopted in the implementation. In other words, neither fuzzy logic nor explicit item multiplicities are used in the thesis, precisely because the objective is to mine frequent explanatory patterns efficiently while preserving the anti-monotonic support structure on which the mining process depends.

Association rules offer several important advantages. They are inherently interpretable, because they can be read directly by domain experts without requiring access to a complex predictive model. They also capture multivariate interactions and support exploratory analysis and subgroup characterization, as highlighted by Hullermeier \cite{hullermeier2006association} and Fournier-Viger et al. \cite{fournierviger2017survey}. At the same time, the literature recognizes important limitations: ARM can generate very large and redundant rule sets, it is sensitive to preprocessing decisions, and naive reliance on confidence alone can be misleading. For this reason, modern ARM depends heavily on post-processing procedures and on multiple interestingness criteria, again in line with the discussions of Hullermeier \cite{hullermeier2006association} and McNicholas and Murphy \cite{mcnicholas2008lift}.

In the proposed pipeline, ARM is not treated as a standalone predictive technique, but as the mechanism through which explanation-derived evidence is synthesized at a global level. Counterfactual analysis provides local evidence about which features must change in order to alter the prediction; ARM then summarizes how these relevant features systematically co-occur across many explanations. In this setting, the extracted rules describe recurring structures in the explanations themselves rather than simple co-occurrence patterns in the raw data. More specifically, the approach leverages Boolean ARM with FP-Growth to efficiently extract frequent explanatory patterns while preserving the anti-monotonic support structure required by classical frequent-itemset mining. This role aligns closely with the objective of the thesis: to exploit the interpretability, multivariate expressive power, and filtering mechanisms of ARM in order to transform local counterfactual instances into a compact and human-readable global description of feature interactions and potential bias regularities.

\section{Literature Gap and Positioning of This Thesis}
The literature reviewed in this chapter provides a solid foundation for each component of the thesis, but it only partially addresses their integration within a unified explanatory pipeline. XAI research offers mature conceptual taxonomies and well-established interpretive objectives; counterfactual-explanation research provides effective local tools for analyzing decision changes; feature-selection literature clarifies how relevance can be formalized; and ARM literature offers efficient mechanisms for extracting recurrent patterns from structured data. What remains less explicitly developed is a framework that starts from local explanations, derives classifier-independent feature relevance from them, and then organizes that relevance into interpretable global patterns.

More specifically, the literature only partially covers the joint connection among:
\begin{itemize}
  \item local counterfactual explanations,
  \item feature relevance assessed independently of the specific classifier and inferred from decision-changing evidence,
  \item global pattern extraction through Boolean association-rule mining,
  \item and the subsequent analysis of rules that are actionable or potentially biased.
\end{itemize}

This is the precise point at which the thesis is positioned. The central idea is not merely to explain a classifier at the local level, but to use counterfactual explanations to identify which features most strongly influence the decision, independently of the particular predictive model adopted. In this sense, feature relevance is derived from the variables that repeatedly emerge in counterfactual changes, and these relevant features become the items used to build the ARM transactions. On top of this representation, the thesis applies Boolean ARM with FP-Growth and evaluates the extracted rules through support, confidence, and lift in order to identify stable global patterns, feature correlations, and explanatory regularities that would remain difficult to observe from isolated local explanations alone.

A second element that characterizes the thesis is the interpretive use of the extracted rules. The objective is not merely to mine frequent associations, but to examine which rules are actionable and which may reveal biased or otherwise problematic regularities. In this way, the framework connects local explanation, global feature relevance, rule-based pattern extraction, and bias-sensitive interpretation within a single methodological chain.

Finally, the thesis is also motivated by the practical difficulties associated with real-world datasets, including outdated information, incompleteness, inconsistent usage, and insufficient normalization. These data-related issues are not discussed in detail in this chapter because they are addressed in the following one, but they are already relevant for understanding the methodological choices made here. The need for robust local explanations, structured feature relevance, and interpretable rule extraction emerges precisely in settings where data quality cannot be taken for granted.

Overall, the state of the art supports the core methodological choice of the thesis: starting from local counterfactual explanations, deriving classifier-independent feature relevance, transforming those relevant features into Boolean ARM transactions, and using FP-Growth to extract frequent patterns that can then be interpreted through actionable and bias-related rules. The following chapters build directly on this rationale by describing the dataset and its criticalities, formalizing the method, and evaluating it empirically.